% !PRINTED 209
\weberpage{209}
Ist nun $S$ irgend eine Substitution der Gruppe $\mathfrak{G}$, durch die, wenn $a, b, c, \partial$ ganze Zahlen bedeuten, die nach dem Modul $n$ genommen sind,

\[
x_{1,0}, x_{0,1} \quad \text{in} \quad x_{\partial,-c}, x_{-b,a}
\]

übergeht, so können wir $S$ auf jede der Formeln (1), (2) anwenden. Es ist aber nach (2):

\[
x_{\mu,0} = \varphi_\mu(x_{1,0}),
\]

woraus zu schließen, daß durch die Substitution $S$

\[
x_{\mu,0} \quad \text{in} \quad \varphi_\mu(x_{\partial,-c}) = x_{\partial\mu,\,-c\mu}
\]

übergeht. Ebenso findet man, daß durch $S$

\[
x_{0,\mu'} \quad \text{in} \quad \varphi_{\mu'}(x_{-b,a}) = x_{-b\mu',\,a\mu'}
\]

übergeht.

Wenden wir dies auf die aus (1) hervorgehende Gleichung

\[
x_{\mu,\mu'} = f(x_{\mu,0}, x_{0,\mu'})
\]

an, so folgt, daß durch $S$

\weq{3}{x_{\mu,\mu'} \quad \text{in} \quad x_{\partial\mu - b\mu',\,-c\mu + a\mu'}}

übergeht.

Hieraus schließen wir zunächst, daß die Zahlen $a, b, c, \partial$ so beschaffen sein müssen, daß ihre Determinante

\weq{4}{m = a\partial - bc}

mit $n$ keinen Teiler gemein hat. Denn wäre ein solcher Teiler vorhanden, so würden $\mu, \mu'$, ohne daß beide durch $n$ teilbar sind, so bestimmt werden können, daß

\[
\partial\mu - b\mu' \equiv 0, \quad -c\mu + a\mu' \equiv 0 \pmod{n},
\]

und es würden mehrere voneinander verschiedene Wurzeln $x_{\mu,\mu'}$ durch $S$ in ein und dieselbe Wurzel übergehen, was unmöglich ist. Denn zunächst können weder $a$ und $b$ noch $c$ und $\partial$ einen Teiler mit $n$ gemein haben, weil sonst, wenn $a\mu'$, $b\mu'$ oder $c\mu$, $\partial\mu$ durch $n$ teilbar genommen werden, alle $x_{\mu',0}$ oder alle $x_{0,\mu}$ in $x_{0,0}$ übergehen würden. Setzt man dann $m\mu = n(ha + h'b)$, $m\mu' = n(hc + h'\partial)$, und bestimmt $h$ und $h'$ so, daß $ha + h'b$ keine Teiler mit $n$ gemein hat, so brauchen, wenn $n$ und $m$ einen gemeinsamen Teiler haben, $\mu, \mu'$ nicht beide durch $n$ teilbar zu sein, und es geht nicht nur $x_{0,0}$, sondern auch $x_{\mu,\mu'}$ nach (3) in $x_{0,0}$ über.

Bezeichnen wir abgekürzt die Vertauschung (3) mit

\[
(x_{\mu,\mu'}, x_{\nu,\nu'}),
\]

so ist

\[
\begin{aligned}
\nu &\equiv \partial\mu - b\mu' \\
\nu' &\equiv -c\mu + a\mu'
\end{aligned}
\pmod{n},
\]
%CONT>>
